\documentclass[conference]{IEEEtran}
\usepackage[utf8]{inputenc}
\usepackage{graphicx}
\usepackage{booktabs}
\usepackage{array}
\usepackage{url}

\title{Proposed Solution Report:\\A Cloud-Based Movie Streaming Platform Using FastAPI and Next.js}

\author{
\IEEEauthorblockN{Mustakim Ahmed Hasan}
\IEEEauthorblockA{
Department of Computer Science and Engineering\\
North South University\\
Email: mustakimahmedhasan@gmail.com
}
}

\begin{document}
\maketitle

\begin{abstract}
This report presents the proposed solution for a full-stack movie streaming platform developed as a CSE299 project. The system provides secure user authentication, movie discovery, adaptive video playback with HTTP Live Streaming (HLS), watch-history tracking, and recommendation endpoints. The architecture uses a decoupled frontend-backend model: Next.js for the client application, FastAPI for REST APIs, PostgreSQL for metadata persistence, and Supabase Storage for media objects. Deployment is planned as cloud-hosted frontend and backend services with environment-based configuration for scalability and maintainability. This document defines the problem, architecture, implementation plan, deployment strategy, known risks, and evaluation metrics.
\end{abstract}

\begin{IEEEkeywords}
FastAPI, Next.js, HLS, PostgreSQL, Supabase Storage, cloud deployment, streaming platform
\end{IEEEkeywords}

\section{Introduction}
Conventional academic web projects often focus on static CRUD workflows and do not cover real media-delivery constraints such as adaptive streaming, secure token-based sessions, and cross-origin deployment. This project targets a practical Netflix-style platform that combines modern frontend engineering, backend API design, and cloud storage integration into one deployable system.

\section{Problem Statement}
Users need a web application where they can register, log in, browse content, search movies, watch streams, and resume progress. Existing classroom prototypes usually fail in one or more of these areas:
\begin{itemize}
\item Incomplete authentication and session handling.
\item Poor UX for playback and content discovery.
\item Missing cloud-ready architecture for media and APIs.
\item Weak separation between frontend and backend concerns.
\end{itemize}

\section{Objectives and Scope}
\subsection{Objectives}
\begin{itemize}
\item Build a responsive streaming frontend for desktop and mobile browsers.
\item Provide secure backend APIs for auth, content, and watch history.
\item Integrate HLS playback with subtitle track support.
\item Persist metadata and user records in PostgreSQL.
\item Store and serve media via cloud object storage.
\item Deploy frontend and backend as independent cloud services.
\end{itemize}

\subsection{In Scope}
\begin{itemize}
\item Authentication: register, login, logout, current-user profile.
\item Movie catalog APIs: list, search, detail, trending, recommendations.
\item Playback page using HLS URLs from backend metadata.
\item Watch-history read and upsert APIs.
\item CORS and environment configuration for cross-domain operation.
\end{itemize}

\subsection{Out of Scope}
\begin{itemize}
\item Payment or subscription billing.
\item DRM-grade content protection.
\item Native mobile applications.
\item Real-time chat/social features.
\end{itemize}

\section{System Architecture}
The solution follows a three-layer architecture:
\begin{itemize}
\item \textbf{Presentation layer:} Next.js (React, TypeScript, SWR, Tailwind CSS, HLS.js).
\item \textbf{Service layer:} FastAPI with modular routers (\texttt{/auth}, \texttt{/movies}, \texttt{/history}).
\item \textbf{Data layer:} PostgreSQL (application metadata) and Supabase Storage (media assets).
\end{itemize}

Frontend and backend communicate through REST endpoints under \texttt{/api/v1}. Cookies are used for authentication continuity, and CORS middleware controls allowed origins.

\section{Technology Stack and Justification}
\begin{table}[h]
\caption{Selected Technologies}
\label{tab:stack}
\centering
\begin{tabular}{p{2.1cm} p{4.8cm}}
\toprule
\textbf{Layer} & \textbf{Technology and Rationale} \\
\midrule
Frontend & Next.js with TypeScript for structured UI development; SWR for data fetching and caching; Tailwind CSS for fast UI iteration; HLS.js for browser-compatible streaming. \\
Backend & FastAPI for async REST API development; SQLAlchemy async ORM for maintainable DB interaction; Uvicorn ASGI server. \\
Database & PostgreSQL for relational consistency across users, movies, genres, subtitles, and watch history. \\
Media Storage & Supabase Storage for cloud-hosted HLS playlists and segments. \\
Deployment & Cloud deployment with environment-based configuration and CORS control. \\
\bottomrule
\end{tabular}
\end{table}

\section{Data Model Overview}
Core entities include \texttt{User}, \texttt{Movie}, \texttt{Genre}, \texttt{MovieGenre}, \texttt{Subtitle}, and \texttt{WatchHistory}. Key relationships:
\begin{itemize}
\item One user to many watch-history records.
\item One movie to many subtitle records.
\item Many-to-many mapping between movies and genres.
\end{itemize}

This model supports filtering, recommendation heuristics, and playback metadata retrieval.

\section{API Design Summary}
Implemented route groups:
\begin{itemize}
\item \textbf{Auth:} \texttt{/auth/register}, \texttt{/auth/login}, \texttt{/auth/logout}, \texttt{/auth/me}
\item \textbf{Movies:} \texttt{/movies}, \texttt{/movies/\{id\}}, \texttt{/movies/trending}, \texttt{/movies/\{id\}/recommendations}
\item \textbf{History:} \texttt{/history} (GET and POST upsert)
\end{itemize}

The backend uses typed response schemas so frontend pages can rely on consistent field names such as \texttt{hls\_master\_url}, \texttt{duration\_minutes}, and \texttt{maturity\_rating}.

\section{Frontend UX Plan}
Primary user flows:
\begin{itemize}
\item Home: featured hero, trending row, latest row.
\item Search: query-driven movie grid.
\item Movie detail: player, metadata, recommendations.
\item History: per-user watched progress list.
\item Auth pages: login and register.
\end{itemize}

SWR is used for asynchronous data loading with state-based rendering for loading, success, and error views.

\section{Deployment Strategy}
\subsection{Frontend}
Deploy Next.js app with environment variable:
\begin{itemize}
\item \texttt{NEXT\_PUBLIC\_API\_BASE=<backend\_url>}
\end{itemize}

\subsection{Backend}
Deploy FastAPI as a cloud web service with:
\begin{itemize}
\item \texttt{DATABASE\_URL} for PostgreSQL (pooler-based in cloud).
\item \texttt{SECRET\_KEY} for token generation.
\item \texttt{CORS\_ORIGINS} as comma-separated or JSON origins list.
\end{itemize}

\subsection{Media}
Store HLS master playlists and segments in Supabase bucket paths (for example: \texttt{media/mandalorian/...}) and persist public URLs in DB movie records.

\section{Implementation Plan and Timeline}
\begin{table}[h]
\caption{Planned Work Breakdown}
\label{tab:timeline}
\centering
\begin{tabular}{p{1.2cm} p{5.7cm}}
\toprule
\textbf{Week} & \textbf{Milestone} \\
\midrule
1--2 & Backend skeleton, models, auth routes, DB session integration. \\
3--4 & Movie APIs, history APIs, recommendation logic, schema validation. \\
5--6 & Next.js pages, player integration, search/history UX, type fixes. \\
7 & Cloud media integration (Supabase Storage), endpoint wiring. \\
8 & Deployment hardening: CORS, env setup, cross-origin testing. \\
9 & Final QA, report preparation, demo packaging. \\
\bottomrule
\end{tabular}
\end{table}

\section{Risk Analysis and Mitigation}
\begin{itemize}
\item \textbf{CORS instability:} normalize origins and verify preflight responses in production.
\item \textbf{Database connectivity in serverless/cloud:} use pooler endpoints and compatible async driver settings.
\item \textbf{Schema drift between frontend and backend:} enforce typed API contracts.
\item \textbf{Large media handling:} keep media outside repository; use object storage/CDN.
\end{itemize}

\section{Evaluation Metrics}
Success will be measured with:
\begin{itemize}
\item Functional coverage of all planned routes and pages.
\item End-to-end playback success from UI to HLS source.
\item Successful auth and watch-history persistence.
\item Stable cloud deployment without recurrent runtime failures.
\item Acceptable page load and response latency for demo use.
\end{itemize}

\section{Expected Outcome}
The final system is expected to provide a realistic, cloud-deployed streaming prototype that demonstrates full-stack engineering competence across API development, data modeling, frontend UX, media delivery, and deployment troubleshooting.

\section{Conclusion}
The proposed solution is feasible within the project timeline and technically aligned with course expectations for applied software engineering. The architecture is modular, deployable, and extensible, making it suitable for both academic demonstration and future enhancement.

\begin{thebibliography}{00}
\bibitem{fastapi} FastAPI Documentation. [Online]. Available: \url{https://fastapi.tiangolo.com}
\bibitem{nextjs} Next.js Documentation. [Online]. Available: \url{https://nextjs.org/docs}
\bibitem{sqlalchemy} SQLAlchemy Documentation. [Online]. Available: \url{https://docs.sqlalchemy.org}
\bibitem{supabase} Supabase Documentation. [Online]. Available: \url{https://supabase.com/docs}
\bibitem{hls} HTTP Live Streaming (HLS) Overview. [Online]. Available: \url{https://developer.apple.com/streaming/}
\end{thebibliography}

\end{document}
